\[ %%%%%%%%%%%%%%%%%%%%%%%%%%%%%%%%%%%%%%%%%%%%%%%%%%%%%%%%%%%%%%%%%%%%%%%%%%%%%%%% \subsection{Multi-Level Physics Constraint Integration} \label{sec:multi_level_constraints} %%%%%%%%%%%%%%%%%%%%%%%%%%%%%%%%%%%%%%%%%%%%%%%%%%%%%%%%%%%%%%%%%%%%%%%%%%%%%%%% A central design principle of PHYSGEN is that physical knowledge should be integrated at multiple hierarchical levels of the computational architecture. Unlike conventional physics-informed approaches, where physical constraints are commonly introduced as additional terms in the optimization objective, PHYSGEN incorporates physical information directly into the representation, interaction, evolution, and stability mechanisms of the model. The proposed multi-level physics constraint framework consists of four complementary levels: \begin{enumerate} \item node-level physical constraints; \item edge-level interaction constraints; \item graph-level structural constraints; \item system-level conservation constraints. \end{enumerate} This hierarchical organization enables PHYSGEN to preserve both local physical consistency and global dynamical behavior. %%%%%%%%%%%%%%%%%%%%%%%%%%%%%%%%%%%%%%%%%%%%%%%%%%%%%%%%%%%%%%%%%%%%%%%%%%%%%%%% \subsubsection{Node-Level Physical Constraints} %%%%%%%%%%%%%%%%%%%%%%%%%%%%%%%%%%%%%%%%%%%%%%%%%%%%%%%%%%%%%%%%%%%%%%%%%%%%%%%% At the lowest level, PHYSGEN constrains individual node representations according to physically meaningful state variables. For a node $v_i$, the latent representation is defined as: \begin{equation} h_i^t=f_{enc}(x_i^t), \end{equation} where $x_i^t$ contains observable physical quantities. The node constraint function is defined as: \begin{equation} \mathcal{C}_{node} = \sum_i \| \Gamma_i(h_i^t)-x_i^t \|^2, \end{equation} where $\Gamma_i$ represents a physical interpretation mapping from latent space to measurable quantities. Node-level constraints enforce: \begin{itemize} \item physically valid state ranges; \item dimensional consistency; \item positivity conditions; \item material or domain-specific restrictions. \end{itemize} For example, temperature, density, concentration, or energy variables must remain within physically admissible domains. %%%%%%%%%%%%%%%%%%%%%%%%%%%%%%%%%%%%%%%%%%%%%%%%%%%%%%%%%%%%%%%%%%%%%%%%%%%%%%%% \subsubsection{Edge-Level Interaction Constraints} %%%%%%%%%%%%%%%%%%%%%%%%%%%%%%%%%%%%%%%%%%%%%%%%%%%%%%%%%%%%%%%%%%%%%%%%%%%%%%%% Physical systems evolve through interactions between components. Therefore, PHYSGEN introduces constraints directly on graph edges. The interaction representation is: \begin{equation} e_{ij}^{t} = f_e(x_i^t,x_j^t,p_{ij}), \end{equation} where $p_{ij}$ contains physical interaction parameters. The edge consistency constraint is defined as: \begin{equation} \mathcal{C}_{edge} = \sum_{(i,j)\in E} \| e_{ij}^{t} - \hat{e}_{ij}^{t} \|^2. \end{equation} This constraint regulates: \begin{itemize} \item local interaction strength; \item spatial coupling; \item directional dependencies; \item physical connectivity. \end{itemize} Therefore, message propagation is guided by physically meaningful relationships rather than purely statistical correlations. %%%%%%%%%%%%%%%%%%%%%%%%%%%%%%%%%%%%%%%%%%%%%%%%%%%%%%%%%%%%%%%%%%%%%%%%%%%%%%%% \subsubsection{Graph-Level Structural Constraints} %%%%%%%%%%%%%%%%%%%%%%%%%%%%%%%%%%%%%%%%%%%%%%%%%%%%%%%%%%%%%%%%%%%%%%%%%%%%%%%% At the graph level, PHYSGEN preserves structural properties of the complete physical system. The graph evolution is represented as: \begin{equation} G_t \rightarrow G_{t+1}. \end{equation} The structural constraint is defined as: \begin{equation} \mathcal{C}_{graph} = D(G_t,G_{t+1}), \end{equation} where $D(\cdot)$ measures deviation between consecutive graph states. Depending on the physical domain, graph-level constraints may preserve: \begin{itemize} \item topology consistency; \item symmetry properties; \item connectivity preservation; \item boundary conditions; \item geometric relationships. \end{itemize} This level is particularly important for systems where spatial organization influences future evolution. %%%%%%%%%%%%%%%%%%%%%%%%%%%%%%%%%%%%%%%%%%%%%%%%%%%%%%%%%%%%%%%%%%%%%%%%%%%%%%%% \subsubsection{System-Level Conservation Constraints} %%%%%%%%%%%%%%%%%%%%%%%%%%%%%%%%%%%%%%%%%%%%%%%%%%%%%%%%%%%%%%%%%%%%%%%%%%%%%%%% At the highest level, PHYSGEN incorporates global physical invariants. Let: \begin{equation} I(X_t) \end{equation} represent a conserved system quantity. The conservation error is: \begin{equation} \mathcal{C}_{sys} = \| I(X_{t+1}) - I(X_t) \|^2. \end{equation} Examples include: \begin{itemize} \item mass conservation; \item energy conservation; \item momentum conservation; \item charge conservation; \item probability conservation. \end{itemize} The system-level constraint prevents physically impossible long-term drift. %%%%%%%%%%%%%%%%%%%%%%%%%%%%%%%%%%%%%%%%%%%%%%%%%%%%%%%%%%%%%%%%%%%%%%%%%%%%%%%% \subsubsection{Unified Multi-Level Physics Objective} %%%%%%%%%%%%%%%%%%%%%%%%%%%%%%%%%%%%%%%%%%%%%%%%%%%%%%%%%%%%%%%%%%%%%%%%%%%%%%%% The complete physics-guided objective of PHYSGEN combines all constraint levels: \begin{equation} \mathcal{L}_{physics} = \lambda_n \mathcal{C}_{node} + \lambda_e \mathcal{C}_{edge} + \lambda_g \mathcal{C}_{graph} + \lambda_s \mathcal{C}_{sys}, \end{equation} where: \begin{itemize} \item $\lambda_n$ controls node-level constraints; \item $\lambda_e$ controls interaction constraints; \item $\lambda_g$ controls graph structure constraints; \item $\lambda_s$ controls global conservation constraints. \end{itemize} The complete PHYSGEN optimization objective becomes: \begin{equation} \mathcal{L}_{PHYSGEN} = \mathcal{L}_{prediction} + \mathcal{L}_{physics} + \lambda_{stab} \mathcal{L}_{stability}. \end{equation} However, unlike traditional physics-informed neural networks, these constraints influence not only parameter optimization but also intermediate computational transformations. %%%%%%%%%%%%%%%%%%%%%%%%%%%%%%%%%%%%%%%%%%%%%%%%%%%%%%%%%%%%%%%%%%%%%%%%%%%%%%%% \subsubsection{Hierarchical Physics Guidance Principle} %%%%%%%%%%%%%%%%%%%%%%%%%%%%%%%%%%%%%%%%%%%%%%%%%%%%%%%%%%%%%%%%%%%%%%%%%%%%%%%% The multi-level constraint architecture establishes a hierarchical physics guidance principle: \begin{equation} Physical\ Knowledge \rightarrow Representation \rightarrow Interaction \rightarrow Evolution \rightarrow Prediction. \end{equation} This principle distinguishes PHYSGEN from approaches where physical laws are considered only external regularizers. By embedding physical constraints throughout the computational hierarchy, PHYSGEN aims to achieve improved generalization, enhanced stability, and physically consistent long-horizon forecasting. %%%%%%%%%%%%%%%%%%%%%%%%%%%%%%%%%%%%%%%%%%%%%%%%%%%%%%%%%%%%%%%%%%%%%%%%%%%%%%%% \subsubsection{Summary} %%%%%%%%%%%%%%%%%%%%%%%%%%%%%%%%%%%%%%%%%%%%%%%%%%%%%%%%%%%%%%%%%%%%%%%%%%%%%%%% The multi-level physics constraint integration mechanism provides a unified framework for incorporating scientific knowledge into neural dynamical models. By simultaneously constraining individual states, interactions, graph structures, and global invariants, PHYSGEN creates a physically guided computational environment where learning and physical reasoning are continuously coupled. This hierarchical integration forms the foundation for reliable scientific prediction and enables PHYSGEN to serve as a general framework for physics-guided dynamical modeling. \] 